\begin{lemma}[The identity map inhabits the quasiregularity data of \S5 at $K=1$]
Let $n\ge2$ be an integer and let $\Omega\subseteq\R^n$ be open.  Put
\[
F(x)=x,\qquad G^i(x)=e_i\quad(i=1,\dots,n),\qquad \tau(x)=I,
\]
where $e_i$ is the $i$-th standard basis vector and $I$ is the $n\times n$
identity matrix.  Then
\begin{enumerate}
\item $F$ is $1$-quasiregular on $\Omega$ with $G^i$ the weak gradient of the
  $i$-th component of $F$ (\noderef{IsQuasiregularOn} at $K=1$); and
\item $\tau$ is the distortion matrix of $(F,G)$ on $\Omega$
  (\noderef{IsDistortionMatrix}).
\end{enumerate}

\noindent\textbf{Why this node exists.}  The Note for the formalization requires
at l.~258 that a cited result ``must not be vacuous (each must be a true
statement for the operators to which it is applied)'', and the demand carries to
the data hypotheses the \S5 statements are conditioned on: Theorem~Theorem~7 of the thesis
assumes \noderef{IsQuasiregularOn} together with \noderef{IsDistortionMatrix},
and if no pair $(F,G,\tau)$ satisfied them the theorem and everything proved from
it would be vacuously true.  The identity map is the paper's own trivial
instance.  The thesis defines $K$-quasiregularity at l.~1582 as
$F\in W^{1,n}_{\mathrm{loc}}(\Omega)$ together with
\[
|DF(x)|^n\le K\,J_F(x)\quad\text{for a.e. }x\in\Omega,
\tag{5.10}
\]
and for $F=\mathrm{id}$ one has $DF=I$, so $|DF(x)|=1$ and $J_F(x)=1$ and
(5.10) of the thesis holds with equality at $K=1$, which is the least admissible
value since Theorem~Theorem~7 of the thesis assumes $K\ge1$.  This node is the companion of the
model-field and model-density witnesses for the
\S5 data, and like them it is re-checked mechanically whenever the predicates it
inhabits are edited.

\textbf{What it does not claim.}  It says nothing about the cited result (C7).
(C7) is a \emph{cited} statement: per the Note's closed list it is imported, not
proved here, so ``non-vacuity'' for it can only mean that its hypotheses are
inhabited --- which is exactly what clauses (1) and (2) above supply, since
(C7)'s hypotheses are $K$-quasiregularity of $F$ and that $\tau$ is its
distortion matrix.  Nor does it claim $F$ is nonconstant-dependent: the identity
map is not constant, so the nonconstancy hypothesis that
Theorem~Theorem~7 of the thesis and the distortion-matrix properties lemma carry is also
available for this witness, but it is not part of either predicate and so is not
asserted here.
\end{lemma}
\begin{proof}
Throughout, $e_i=$ the $i$-th standard basis vector, so that its $j$-th
coordinate is $e_i^j=\delta_{ij}$, and the matrix whose $(i,j)$ entry is
$G^i(x)^j=\delta_{ij}$ is the identity matrix $I$, independently of $x$.

\emph{Part 1(a): each component of $F$ lies in $W^{1,n}_{\mathrm{loc}}(\Omega)$
with weak gradient $e_i$.}  Fix $i$ and let $u(x)=F(x)^i=x^i$, the $i$-th
coordinate function, and let $g(x)=e_i$ for all $x$.  We verify the hypotheses of
\noderef{WeakGradientOfC1} with $q=n$:
\begin{itemize}
\item $n\ge1$ and $q=n\ge2\ge1$, using $n\ge2$;
\item $\Omega$ is open, by hypothesis;
\item $u$ is a linear functional on the finite-dimensional space $\R^n$, hence
  differentiable at every point of $\R^n$ --- in particular at every point of
  $\Omega$ --- with $Du(x)=u$ for every $x$; therefore, for every $x$ and every
  index $j$,
  \[
  D_ju(x)=Du(x)[e_j]=u(e_j)=e_j^i=\delta_{ij}=e_i^j=g(x)^j ,
  \]
  which is the required identity $D_ju(x)=g^j(x)$ for all $x\in\Omega$ and all
  $j$;
\item $g$ is constant, hence continuous on $\Omega$.
\end{itemize}
So \noderef{WeakGradientOfC1} gives
$u\in W^{1,n}_{\mathrm{loc}}(\Omega)$ with weak gradient $g$, i.e.\ the $i$-th
component of $F$ lies in $W^{1,n}_{\mathrm{loc}}(\Omega)$ with weak gradient
$G^i$.  As $i$ was arbitrary, the first clause of \noderef{IsQuasiregularOn}
holds.

\emph{Part 1(b): the pointwise distortion inequality at $K=1$.}  Let
$x\in\Omega$ be arbitrary and let $\xi\in\R^n$.  Since $G^i(x)^j=\delta_{ij}$,
\[
\sum_{j=1}^nG^i(x)^j\xi^j=\sum_{j=1}^n\delta_{ij}\xi^j=\xi^i ,
\]
so, using $\sum_i(\xi^i)^2=|\xi|^2$ for the Euclidean norm,
\[
\Bigl(\sum_{i=1}^n\Bigl(\sum_{j=1}^nG^i(x)^j\xi^j\Bigr)^{\!2}\Bigr)^{\!n/2}
=\bigl(|\xi|^2\bigr)^{n/2}=|\xi|^{\,n},
\]
the last equality because $|\xi|\ge0$ and $(t^2)^{n/2}=t^n$ for $t\ge0$.  The
matrix $\bigl(G^i(x)^j\bigr)_{i,j}$ is $I$, whose determinant is $1$.  Hence the
right-hand side of the inequality in \noderef{IsQuasiregularOn} at $K=1$ is
\[
1\cdot\det\bigl(G^i(x)^j\bigr)_{i,j}\cdot|\xi|^{\,n}=1\cdot1\cdot|\xi|^{\,n}
=|\xi|^{\,n},
\]
so the inequality holds with equality.  Since $x\in\Omega$ and $\xi$ were
arbitrary, it holds at \emph{every} point of $\Omega$ and for every $\xi$, hence
for almost every $x$ with $x\in\Omega$, which is the second clause of
\noderef{IsQuasiregularOn}.  This proves (1).

\emph{Part 2: $\tau=I$ is the distortion matrix.}  Let $x\in\Omega$.  By the
first paragraph the matrix $M=\bigl(G^i(x)^j\bigr)_{i,j}$ is $I$, so
$\det M=1$.  \noderef{IsDistortionMatrix} asks for two implications at $x$.

For the first, its hypothesis $\det M>0$ holds, since $\det M=1>0$, and we must
check
\[
\tau(x)\,\bigl(M^{\mathsf T}M\bigr)=(\det M)^{2/n}\,I .
\]
Here $M=I$ gives $M^{\mathsf T}=I$ and $M^{\mathsf T}M=I$, and $\tau(x)=I$, so
the left-hand side is $I\cdot I=I$.  The right-hand side is
$1^{2/n}\,I=1\cdot I=I$, because $1^t=1$ for every real $t$.  The two sides agree.

For the second implication, its hypothesis is $\neg(\det M>0)$, which is false
because $\det M=1>0$; so the implication holds vacuously.  (Had it been invoked,
its conclusion $\tau(x)=I$ is true here as well.)

As $x\in\Omega$ was arbitrary, $\tau$ is the distortion matrix of $(F,G)$ on
$\Omega$, which is (2).
\end{proof}
